% Created 2026-06-28 Sun 10:29
% Intended LaTeX compiler: pdflatex
\documentclass[11pt]{article}
\usepackage[utf8]{inputenc}
\usepackage[T1]{fontenc}
\usepackage{graphicx}
\usepackage{longtable}
\usepackage{wrapfig}
\usepackage{rotating}
\usepackage[normalem]{ulem}
\usepackage{amsmath}
\usepackage{amssymb}
\usepackage{capt-of}
\usepackage{hyperref}
\date{\today}
\title{README.org - Minha configuração do Emacs}
\hypersetup{
 pdfauthor={},
 pdftitle={README.org - Minha configuração do Emacs},
 pdfkeywords={},
 pdfsubject={},
 pdfcreator={Emacs 30.2 (Org mode 9.7.11)}, 
 pdflang={English}}
\begin{document}

\maketitle
\tableofcontents

\section{Funções custom}
\label{sec:org9898316}
\subsection{Markdown link to Org link}
\label{sec:org73bd5f6}
O Zen Browser me dá os links em Markdown, então eu preciso converter para Org.
\begin{verbatim}
(defun md-link-to-org-link ()
  "Converte link de Markdown para Org"
  (interactive)
  (when (region-active-p)
    (save-excursion
      (replace-regexp "\\[\\(.*\\)\\](\\(.*\\))" "[[\\2][\\1]]" nil (region-beginning) (region-end))
      )))
(add-hook 'org-mode-hook
          (lambda ()
            (keymap-set org-mode-map "C-c C-g" 'md-link-to-org-link)
            )
          )
\end{verbatim}
\texttt{C-c C-g} agora converte o link markdown selecionado pro formato Org.
\texttt{(interactive)} é necessário para poder fazer o \emph{keybind}.
\texttt{(save-excursion)} grava onde o cursor estava quando o código foi executado e retorna quando acaba.
\texttt{region-active-p} retorna \texttt{t} se tem texto selecionado e \texttt{nil} se não.
\section{Terminal}
\label{sec:orgdf55e80}
Usar terminal no Emacs tem sido a pior coisa dessa experiência de longe. Vamos ver como dá pra melhorar.
\href{https://www.masteringemacs.org/article/running-shells-in-emacs-overview}{Running Shells and Terminal Emulators in Emacs - Mastering Emacs}
\begin{itemize}
\item Todos os shells derivam de um modo chamado \texttt{comint-mode}.
\end{itemize}

E eu fiquei nessa. É inacreditável como é confuso isso. Vou tentar postergar ao máximo ter que usar um terminal.
\end{document}
